% =============================================================================
% CompetitiveExam/CS301/05-queues-answers.tex
% Answer key for 05-queues-questions.tex.
%
% Real GATE PYQs (Q1--Q5) are sourced from the local GATEOverflow Volume 2
% PYQ book (CompetitiveExam/Books/index.md, Tier-1 availability: Queue
% cluster PDF pp.240-244, Priority Queue cluster p.239), which
% authenticates each question (year + question number printed on the
% page). The book's own answer keys are collapsed and not extractable,
% so each real-PYQ answer was independently established in the build
% session (2026-09-17): GATEOverflow online question pages, ExamSIDE
% year-wise/topic listings, Testbook's GATE-2012 solved page,
% appliedroots/SamagraCS/BYJUS solved listings, the GFG GATE-IT-2007
% quiz, the documented monotonic-key principle for stack-from-priority-
% queue, and — where the answer is exact algorithmics — deterministic
% re-derivation. Original questions (Q6--Q8) have answers derived
% directly from the Week 5 deck algorithms, confirmed by execution
% (2026-09-17).
% =============================================================================

\documentclass[11pt]{article}
\usepackage[margin=1in]{geometry}
% =============================================================================
% Preambles/header.tex — shared LaTeX/Beamer preamble for this project
%
% Use from a Beamer lecture by prepending:
%     \input{header}
% and compiling with TEXINPUTS=../Preambles:$TEXINPUTS (the /compile-latex
% skill does this automatically).
%
% PALETTE CONTRACT. Color names defined here MUST match the names in
% Quarto/theme-template.scss so Beamer and Quarto renderings use the same
% palette. The scripts/check-palette-sync.sh script enforces this.
%
% To customize: edit the HEX values below AND the matching $variable in
% Quarto/theme-template.scss, then run scripts/check-palette-sync.sh.
% =============================================================================

% -----------------------------------------------------------------------------
% Engine detection + encoding
%
% Our /compile-latex skill uses XeLaTeX, where inputenc is unnecessary and
% can produce encoding warnings. Only load inputenc under pdfTeX.
% -----------------------------------------------------------------------------
\usepackage{iftex}
\ifPDFTeX
  \usepackage[utf8]{inputenc}
\fi

\usepackage{amsmath, amssymb, amsthm}
\usepackage{booktabs}
\usepackage{graphicx}

% -----------------------------------------------------------------------------
% PALETTE — mirrors Quarto/theme-template.scss variable names.
% Load xcolor and define colors BEFORE anything that references them
% (notably hyperref, hypersetup, and the Beamer color assignments).
% -----------------------------------------------------------------------------
\usepackage{xcolor}

\definecolor{primary-blue}{HTML}{012169}      % headings, accents
\definecolor{primary-gold}{HTML}{B9975B}      % emphasis, borders
\definecolor{highlight-yellow}{HTML}{F2A900}  % markers, alerts
\definecolor{light-bg}{HTML}{E8EDF5}          % title / section backgrounds
\definecolor{jet}{HTML}{1A1A1A}               % body text

% Semantic colors (used in diagrams and annotations)
\definecolor{positive}{HTML}{15803D}          % good / observed / identified
\definecolor{negative}{HTML}{B91C1C}          % bad / problematic / bias
\definecolor{neutral}{HTML}{525252}           % reference / context

% Highlight accents (match .hi-* classes in the SCSS)
\definecolor{hi-slate}{HTML}{314F4F}
\definecolor{hi-green}{HTML}{2E8B57}
\definecolor{hi-red}{HTML}{C41E3A}

% Semantic aliases so skill templates and snippets can write
% \textcolor{accent}{...} without hard-coding "primary-blue".
\colorlet{accent}{primary-blue}
\colorlet{accent2}{primary-gold}

% -----------------------------------------------------------------------------
% Hyperref — loaded AFTER xcolor + palette so link colors resolve.
% -----------------------------------------------------------------------------
\usepackage{hyperref}
\hypersetup{colorlinks=true, linkcolor=primary-blue, urlcolor=primary-blue,
            citecolor=primary-blue, pdfborder={0 0 0}}

% -----------------------------------------------------------------------------
% Course code — set once per deck (after \input{header}, before \begin{document})
% via \coursecode{CS401}. Shown in the footer on every slide so a deck's course
% is identifiable without opening the title page. Empty by default.
% -----------------------------------------------------------------------------
\makeatletter
\newcommand{\coursecode}[1]{\gdef\@coursecodetext{#1}}
\gdef\@coursecodetext{}
\makeatother

% -----------------------------------------------------------------------------
% Beamer theme assignments (only apply under Beamer).
%
% We detect Beamer via \@ifclassloaded{beamer}, which is the documented,
% non-internal way. This must be wrapped in \makeatletter/\makeatother
% because \@ifclassloaded contains an @-catcode character.
% -----------------------------------------------------------------------------
\makeatletter
\@ifclassloaded{beamer}{%
  \usefonttheme{professionalfonts}
  \setbeamercolor{structure}{fg=primary-blue}
  \setbeamercolor{frametitle}{fg=primary-blue}
  \setbeamercolor{title}{fg=primary-blue}
  \setbeamercolor{subtitle}{fg=primary-gold}
  \setbeamercolor{author}{fg=primary-blue}
  \setbeamercolor{institute}{fg=neutral}
  \setbeamercolor{date}{fg=primary-gold}
  \setbeamercolor{itemize item}{fg=primary-blue}
  \setbeamercolor{itemize subitem}{fg=primary-gold}
  \setbeamercolor{alerted text}{fg=primary-gold}
  \setbeamercolor{block title}{fg=white, bg=primary-blue}
  \setbeamercolor{block body}{bg=light-bg}
  % Semantic block variants: block=definition/concept (blue), exampleblock=
  % worked example (green/positive), alertblock=key takeaway or caution (gold).
  % Keep to <=2 colored boxes per slide across ALL three variants (INV-7).
  \setbeamercolor{block title example}{fg=white, bg=positive}
  \setbeamercolor{block body example}{bg=light-bg}
  \setbeamercolor{block title alerted}{fg=white, bg=primary-gold}
  \setbeamercolor{block body alerted}{bg=light-bg}

  % Minimal footer: course code (left) + page X / N (right), no navigation clutter
  \setbeamertemplate{navigation symbols}{}
  \setbeamertemplate{footline}{%
    \color{neutral}\footnotesize\hspace{0.7em}\@coursecodetext\hfill
    \insertframenumber/\inserttotalframenumber\hspace{0.7em}\vspace{0.3em}}
}{}
\makeatother

% -----------------------------------------------------------------------------
% TikZ setup — libraries the snippet gallery uses
% -----------------------------------------------------------------------------
\usepackage{tikz}
\usetikzlibrary{arrows.meta, positioning, calc, decorations.pathreplacing,
                fit, shapes.geometric, backgrounds}

% Shared TikZ styles — snippets and hand-written diagrams can pick these up
% via `\begin{tikzpicture}[every node/.style={...}]` or by naming specific
% styles (e.g., `\node[dag-node] (X) at (0,0) {X};`).
\tikzset{
  % Node styles for causal DAGs and similar
  dag-node/.style = {
    draw=primary-blue, fill=light-bg, rounded corners=2pt,
    minimum width=1.2cm, minimum height=0.9cm, align=center, font=\small
  },
  decision-node/.style = {
    draw=primary-gold, fill=white, diamond, aspect=1.8,
    minimum width=1.8cm, minimum height=1.2cm, align=center, font=\small,
    inner sep=1pt
  },
  % Edge styles — observed vs counterfactual
  observed-edge/.style = {-{Latex[length=2.5mm]}, thick, draw=jet},
  counterfactual-edge/.style = {-{Latex[length=2.5mm]}, thick, draw=neutral, dashed},
  confound-edge/.style = {-{Latex[length=2.5mm]}, thick, draw=negative, dashed},
  % Data-point styles for plots
  observed-dot/.style = {fill=primary-blue, draw=primary-blue, circle, inner sep=1.2pt},
  counterfactual-dot/.style = {fill=white, draw=neutral, circle, inner sep=1.2pt}
}

% -----------------------------------------------------------------------------
% Convenience macros
% -----------------------------------------------------------------------------
\newcommand{\muted}[1]{\textcolor{neutral}{#1}}
\newcommand{\key}[1]{\textcolor{primary-gold}{\textbf{#1}}}
\newcommand{\good}[1]{\textcolor{positive}{#1}}
\newcommand{\bad}[1]{\textcolor{negative}{#1}}

% Standout frame macro: full-bleed dark background for section transitions.
% Beamer-only (uses \begin{frame}/\setbeamercolor/\usebeamerfont) -- do not
% call from an article-class file (e.g. Notes/<CODE>/*-notes.tex). \muted,
% \key, \good, \bad, and \coursecode above are plain xcolor/gdef and are
% safe to use from either class; verified when Notes/article-class support
% was added.
% Expand: \transitionslide{Your transition title}
\newcommand{\transitionslide}[1]{%
  \begingroup
  \setbeamercolor{background canvas}{bg=primary-blue}
  \begin{frame}[plain]
    \centering
    \vfill
    {\usebeamerfont{title}\color{white}\Large #1\par}
    \vfill
  \end{frame}
  \endgroup
}

% Section divider frame (Beamer-only), an alternative to \transitionslide with a
% darker two-line style: near-black (jet) background, a small white label line
% above a larger gold title, and a thin gold rule along the bottom edge.
% Expand: \sectiondivider{Section 2}{Pointers}
\newcommand{\sectiondivider}[2]{%
  \begingroup
  \setbeamercolor{background canvas}{bg=jet}
  \begin{frame}[plain]
    \vspace*{3.0cm}
    {\color{white}\ttfamily\large #1\par}
    \vspace{0.5cm}
    {\color{primary-gold}\LARGE #2\par}
    \begin{tikzpicture}[remember picture, overlay]
      \draw[primary-gold, line width=2.5pt]
        ($(current page.south west)+(0,0.1)$) -- ($(current page.south east)+(0,0.1)$);
    \end{tikzpicture}
  \end{frame}
  \endgroup
}

% -----------------------------------------------------------------------------
% Channel watermark — "Concepts That Click" (YouTube).
%
% Article-class files (CompetitiveExam Q/A sets, Notes, Assignments) get a
% light diagonal draft watermark on every page plus a centered footer naming
% the channel. Beamer decks get a subtle TikZ background watermark instead
% (the footline already carries the course code). Centralized here so every
% MCQ PDF is branded without per-file edits.
% -----------------------------------------------------------------------------
\makeatletter
\@ifclassloaded{article}{%
  \usepackage{draftwatermark}
  \SetWatermarkText{Concepts That Click}
  \SetWatermarkScale{1}
  \SetWatermarkFontSize{2cm}
  \SetWatermarkLightness{0.86}
  \SetWatermarkAngle{45}
  \usepackage{fancyhdr}
  \pagestyle{fancy}
  \fancyhf{}
  \fancyfoot[C]{\footnotesize\textcolor{neutral}{Concepts That Click~|~YouTube: \href{https://www.youtube.com/@concepts.thatclick}{@concepts.thatclick}}}
  \renewcommand{\headrulewidth}{0pt}
  \renewcommand{\footrulewidth}{0pt}
  \fancypagestyle{plain}{%
    \fancyhf{}%
    \fancyfoot[C]{\footnotesize\textcolor{neutral}{Concepts That Click~|~YouTube: \href{https://www.youtube.com/@concepts.thatclick}{@concepts.thatclick}}}%
    \renewcommand{\headrulewidth}{0pt}%
    \renewcommand{\footrulewidth}{0pt}%
  }
}{}
\@ifclassloaded{beamer}{%
  \setbeamertemplate{background}{%
    \begin{tikzpicture}[remember picture, overlay]
      \node[rotate=45, opacity=0.055, align=center]
        at (current page.center)
        {\Huge\textcolor{neutral}{Concepts That Click}};
    \end{tikzpicture}%
  }
}{}
\makeatother 
\coursecode{CS301}
\renewcommand{\thesection}{5.\arabic{section}}

\title{Competitive Exam Practice 5 --- Answer Key\\\large Queues, Circular Buffers, Deques, and Priority Queues}
\author{Auqib Hamid Lone\\[0.3em]
  \emph{Department of Computer Science \& Engineering}, \emph{GCET Ganderbal}}
\date{\today}

\begin{document}
\maketitle

\begin{enumerate}
\item[\textbf{Q1}] \textbf{(A) (ii) and (iii) are true.} (i) is false: stacks are LIFO, so FIFO computations are \emph{not} efficiently supported by them. (ii) is true: linked lists give $O(1)$ insert/delete at a known position versus $O(n)$ shifting on arrays for almost all basic LIST operations. (iii) is true: wraparound reuse keeps every operation $O(1)$ worst case, while the linear two-index version strands free slots (false full) unless it pays $O(n)$ shifting. (iv) is false: queues are FIFO, so LIFO computations are \emph{not} efficiently supported by them. Stem from GATEOverflow V2 p.240; options and answer (a) confirmed against the BYJUS GATE bank and ExamSIDE's GATE-1996 listing (2026-09-17).

\item[\textbf{Q2}] \textbf{(A) \texttt{full}: (\texttt{REAR}+1) mod $n$ == \texttt{FRONT}; \texttt{empty}: \texttt{REAR} == \texttt{FRONT}.} One slot is sacrificed so the two states are distinguishable: with $n-1$ live items \texttt{REAR} sits one slot behind \texttt{FRONT} (full test fires), while the untouched initial state \texttt{REAR} = \texttt{FRONT} = 0 means empty. (B) fails at the start (\texttt{FRONT}+1 $\ne$ \texttt{REAR} on an empty queue); (C) swaps the two tests (initial state would read ``full''); (D) misplaces the $+1$ on \texttt{FRONT}. Confirmed against the GATEOverflow page (accepted reasoning), Testbook's GATE-2012 solution (``option 1''), and ExamSIDE's GATE-2012 listing (2026-09-17).

\item[\textbf{Q3}] \textbf{(A) Both operations can be performed in $O(1)$ time.} A circular-array queue does one array access plus one modular advance per operation --- $O(1)$ worst case for both ENQUEUE and DEQUEUE, with no shifting ever. (B) describes the \emph{linear}-array failure mode (one end drifts into $O(n)$ shifting), not the efficient implementation the stem asks for; (C)/(D) overstate both bills. Stem and options confirmed on appliedroots' GATE-2016-Set-1 listing, the SamagraCS solved page, and ExamSIDE's topic listing; answer (A) per the official GATE 2016 key and the deck's circular implementation (2026-09-17).

\item[\textbf{Q4}] \textbf{(D) strictly decreasing order.} \texttt{DELETEMIN} always returns the smallest key present, but a stack must return the \emph{most recently} pushed item --- so each new key must be smaller than every key already in the queue, i.e.\ strictly decreasing. Increasing/non-decreasing would emulate a queue (FIFO), and non-increasing admits equal keys, which lets an older item with the same key come out first. Tier-1 stem and options from GATEOverflow V2 p.239, year listed on ExamSIDE under GATE CSE 1997, principle cross-checked against the documented monotonic-key construction for stack-from-priority-queue (2026-09-17).

\item[\textbf{Q5}] \textbf{(B) Reverses the order of the elements in the queue.} Each frame deletes the front element, recurses on the rest, then inserts its element at the rear on the way back out --- so the first-deleted (front) element is re-inserted last and ends up at the rear: the queue comes back reversed. (A) would need insert-before-recurse; (C) describes a single rotate, not a full recursion; (D) would need the inserts to be dropped. Confirmed on the GATEOverflow page (best answer B, 56 votes), the GFG GATE-IT-2007 quiz, and ExamSIDE's listing (2026-09-17).

\item[\textbf{Q6}] \textbf{(A) $(5, 3, 5)$.} \texttt{dequeue}: \texttt{front} $= (4+1) \bmod 6 = 5$, \texttt{count} $= 4$ (live items now at indices 5, 0, 1, 2). \texttt{enqueue}: \texttt{rear} $= (2+1) \bmod 6 = 3$, \texttt{count} $= 5$ (items at 5, 0, 1, 2, 3). Not full ($5 < 6$). (B) miscounts \texttt{count} (a dequeue and an enqueue net to zero); (C) forgets to advance \texttt{front}; (D) forgets to advance \texttt{rear}. Confirmed by execution (2026-09-17).

\item[\textbf{Q7}] \textbf{(B) Output-restricted deque; \texttt{dequeueRear} is forbidden.} The restricted cousin that closes the \emph{deletion} end is the output-restricted deque (insert either end, remove one end only) --- here only front deletions are allowed, so rear deletion is the forbidden operation. (A)/(D) misname the variant (input-restricted closes an \emph{insertion} end, and both insertions are open here); (C) names the wrong variant even though it points at the right forbidden end. Cross-checked against the deck's restricted-cousin definitions and the BU Lecture-5 operation sets (2026-09-17).

\item[\textbf{Q8}] \textbf{(A) Unordered array.} Price the workload: unordered pays $\approx 10{,}000 \times O(1)$ inserts plus $5 \times O(n)$ scans ($\sim 60{,}000$ elementary steps); ordered pays $\approx 10{,}000 \times O(n)$ shifting inserts ($\sim 10^8$) to buy $5$ cheap removals --- paying the linear bill on the hot operation. (B)/(D) put the linear cost where the volume is; (C) ignores that the two bills scale with different operation counts. Confirmed by execution (2026-09-17).
\end{enumerate}

\end{document}
